% 原文 pp.285–291（印刷 pp.1403–1409）
% 习题（EXERCISES）。主文件已用 \section*{习题} 起头，故本文件不写 \section，直接从 \subsection* 起。
% 原文每题内部的公式独立编号，依次为 (1)–(10)；此处照录原编号，标签统一取 eq:ex-N。
% 本页范围内原文无插图（无 Figure/示意图），故不含 figure 环境。

\begin{keytakeaways}
本部分是全书末尾的 11 道习题，几乎逐讲对应正文的核心工具。
第 1 题：Goldstone 玻色子与手征微扰；第 2 题：电弱规范理论的谱与两个矢量质量的关系；
第 3 题：超对称 $\sigma$ 模型的大 $N$ 极限；第 4、5 题：Coleman–Mandula 定理所「不允许」的守恒荷与 Lagrange 乘子场；
第 6 题：二维 $U(1)$ 规范理论与 sine–Gordon 模型的 $\theta$ 依赖；第 7 题：$SO(n)$ 规范理论的双值 $\theta$ 角；
第 8 题：$(\psi_+\psi_-)^2$ 模型的无穷多守恒律及 $n=3,4$ 的特殊性；第 9 题：$\mathbb{Z}_n$ orbifold 的基态与射影超曲面上同调；
第 10、11 题：二维 $\mathcal{N}=2$ 理论的两种拓扑扭曲。
解题的公共套路是：先确定对称性破缺或约束之后的物理自由度，再用大 $N$ 鞍点、指标定理或上同调语言读出谱与守恒荷。
\end{keytakeaways}

% ===== 原文 p.285 =====
% （原文页眉：EXERCISES）

\subsection*{Problem 1}

强相互作用在一个很好的近似下由这样一个理论描述：它的对称群 $G=SU(2)\times SU(2)$ 自发破缺（spontaneous symmetry breaking）到对角子群 $H\cong SU(2)$。除 Goldstone 玻色子之外没有无质量粒子；而只要 $G$ 是有效的对称性，这些 Goldstone 玻色子就严格无质量。在很好的近似下，该理论具有一个 $G$ 不变的拉格朗日量（Lagrangian）$L$。

然而，$G$ 并不是自然的严格对称性。它被一些微扰所破坏，这些微扰在 $G$ 的 $(2,2)$ 表示中变换（即两个 $SU(2)$ 各自二维表示的张量积）。在 $G$ 不变理论的算符之中，有两个「多重态」（multiplet）的算符 $\mathcal{O}_i$ 与 $\widetilde{\mathcal{O}}_j$，$i,j=1,\dots,4$，它们各自按 $(2,2)$ 变换。更准确地说，拉格朗日量为

\begin{equation}
  \tag{1}
  \label{eq:ex-1}
  \widetilde{L}=L+\sum_{i=1}^{4}\left(\epsilon_i\mathcal{O}_i+\widetilde{\epsilon}_i\widetilde{\mathcal{O}}_i\right).
\end{equation}

% [Tutor Note: (1) 中两个求和共用同一个哑指标 $i$，但 $\epsilon_i$ 与 $\widetilde{\epsilon}_i$ 属于两个彼此独立的微扰多重态，切勿把 $\widetilde{\epsilon}_i\widetilde{\mathcal{O}}_i$ 误读成 $\epsilon_i$ 的某个函数——作者紧接着特意改用 $\widetilde{\epsilon}_j$ 来提醒这一点。其次，$\mathcal{O}_i$ 与 $\widetilde{\mathcal{O}}_i$ 各有四个分量、均按 $(2,2)$ 变换，故 $\epsilon_i$、$\widetilde{\epsilon}_i$ 本身也是 $(2,2)$ 的四维矢量；只有把 $\epsilon$ 的指标与算符的指标按 $(2,2)$ 配对缩并，$\widetilde{L}$ 才是 $G$ 不变量。物理上 $\epsilon_i$ 编码夸克质量矩阵、$\widetilde{\epsilon}_i$ 编码电磁修正，「恰有两个多重态」这一计数须由底层规范理论（QCD 的手征微扰）给出，本题只把它当作假设。]

$\epsilon_i$ 与 $\widetilde{\epsilon}_j$ 当然全都是实的。（微扰恰好只有 $\mathcal{O}$ 与 $\widetilde{\mathcal{O}}$ 这两个可能的多重态，此事由底层的规范理论解释；但在本题中我们干脆假定它。）

对 $\epsilon_i$ 与 $\widetilde{\epsilon}_i$ 的一般取法，模型的确切对称群只有 $H'=U(1)$；而事实上，这正是自然界的情形。

把有效理论处理到 $\epsilon$ 与 $\widetilde{\epsilon}$ 的一阶（在自然界中它们实际上很小，且彼此大小相当）。

(a) 证明：在一阶，你会预期真空的简并度被解除，并出现一个唯一的、有质量间隙（mass gap）的真空。特别地，你预期 $H'$ 对称性不被破缺。（抽象地说，$H'$ 可以嵌入 $H$ 之中；稍微不平凡的一点是：微扰所选中的那个未微扰理论的真空，在一个确实包含 $H'$ 的 $H\cong SU(2)$ 之下不变。$H$ 当然也可能被共轭成一个不包含 $H'$ 的群。）特别地，不存在被自发破缺的严格对称性，也不存在严格的 Goldstone 玻色子；原先的 Goldstone 玻色子全都获得质量。

(b) 更令人惊讶的是，证明：在 $\epsilon$ 与 $\widetilde{\epsilon}$ 的一阶，三个 Goldstone 玻色子——在 $H'$ 下不变的那一个，以及在 $H'$ 的非平凡表示中变换的那两个——都获得\emph{相等}的质量。（自然界的观察在高精度上也如此：这些粒子分别是中性 $\pi$ 介子与带电 $\pi$ 介子。）

% [Tutor Note: (b) 的「惊讶」在于：一般的 $\epsilon_i$ 取法会破坏同位旋 $SU(2)_V$，按对称性论证三重态的三条质量本应分裂。一阶仍然等质量，是因为三个 Goldstone 模上的质量矩阵正比于单位矩阵：$\mathcal{O}_i$ 的矩阵元只通过 $H'$ 不变的组合进入，而该组合在三重态上各向同性（这正是 Dashen 关于 $\pi$ 介子质量关系的工作所刻画的机制，参见 R. Dashen, Phys. Rev. 183 (1969) 1245）。换言之，质量分裂只能来自 $(\epsilon,\widetilde\epsilon)$ 的二阶，这正是 (c) 所问的阶数；做题时若在一阶就写出分裂，几乎必然是漏掉了「$\mathcal{O}_i$ 按 $(2,2)$ 整体变换」这一约束。]

% ===== 原文 p.286 =====

(c) 三个 Goldstone 玻色子的质量不会\emph{严格}相等。简并度将在 $\epsilon$ 与 $\widetilde{\epsilon}$ 的什么阶被解除？你又如何在有效拉格朗日量中描述这一解除？

\subsection*{Problem 2}

这里我们考虑（如同在弱相互作用的标准模型中那样）四维中的一个 $SU(2)\times U(1)$ 规范理论（gauge theory）。$SU(2)$ 与 $U(1)$ 的规范场将分别称为 $A_2$ 与 $A_1$，其场强（curvature）为 $F_2$ 与 $F_1$，规范耦合（gauge coupling）为 $e_2$ 与 $e_1$。纯规范理论的拉格朗日量将是

\begin{equation}
  \tag{2}
  \label{eq:ex-2}
  L=\int d^4x\left(\frac{1}{4e_1^2}|F_1|^2+\frac{1}{4e_2^2}\operatorname{Tr}F_2^2\right),
\end{equation}

其中 $\operatorname{Tr}$ 是在 $SU(2)$ 的二维表示中取的迹。

% [Tutor Note: 归一化是本式最易出错处。$|F_1|^2$ 是 $F_{1\mu\nu}F_1^{\mu\nu}$（对时空指标缩并，不带 $1/2$）；而 $\operatorname{Tr}F_2^2$ 中的迹若取 $\operatorname{Tr}(T^aT^b)=\frac12\delta^{ab}$，则 $\operatorname{Tr}F_2^2=\frac12F^a_{\mu\nu}F^{a\,\mu\nu}$，两者相差因子 $2$。这个因子直接决定 $e_1,e_2$ 与标准模型惯用记号 $g,g'$ 的换算，并进而影响 (c) 中 $m_W/m_Z=\cos\theta_W$ 的数值；若把两种约定混用，会得到 $\rho$ 参数差 $2$ 的错误结果。另外 (3) 里的 $\lambda/8$ 也不是常见的 $\lambda/4$，同样属约定问题，做题时应从头把 $|\phi|^2=v^2$ 处的二次项系数算一遍而不要照搬标准模型公式。]

此外还有复标量场 $\phi$，它们在 $SU(2)$ 的二维表示中变换，并且对 $U(1)$ 带「电荷 $1/2$」。它们的拉格朗日量为

\begin{equation}
  \tag{3}
  \label{eq:ex-3}
  L_\phi=\int d^4x\left(|D\phi|^2+\frac{\lambda}{8}\left(|\phi|^2-v^2\right)^2\right).
\end{equation}

$v^2$ 是正的。

(a) 证明 $SU(2)\times U(1)$ 自发破缺到一个对角的 $U(1)$。

(b) 确定模型的谱：哪些是无质量的、哪些是有质量的标量与矢量，它们在未破缺的 $U(1)$ 下如何变换？

(c) 模型中观测精度最高的量是两个规范耦合，以及带电与中性有质量矢量的质量。这就给出四个可观测的量。由于模型有四个参数（$e_1$、$e_2$、$\lambda$ 与 $v$），人们也许不会指望这四个可观测的量之间存在任何关系。然而，在经典近似下（由于模型是弱耦合的，这是一个好的近似），这四个可观测的量都与 $\lambda$ 无关。所以在实际上，模型确实预言了两个矢量质量之间的一个关系（在 $e_1$ 与 $e_2$ 已测得的前提下）。这个关系是电弱标准模型早期的成功之一。

当然，这个关系不会严格成立。它会被圈修正（loop correction）所修改，而圈修正将依赖于自然界中存在、但迄今未曾提到的其他场（它们也不影响这里所作的经典讨论）。目前的实验精度刚好足以看到圈效应，从而对参数 $\lambda$ 给出一个估计——$\lambda$ 不进入经典关系，却进入圈修正。

\subsection*{Problem 3}

考虑从 $\mathbb{R}^{2,2}$ 到 $S^N$ 的映照（map）的超对称模型。我们过去曾对映照 $X:\mathbb{R}^{2,2}\to M$（$M$ 为任意 Riemann 流形）写出过拉格朗日量：

\begin{equation}
  \tag{4}
  \label{eq:ex-4}
  L=\int d^2x\,d^2\theta\; g_{IJ}(X)D_+X^I D_-X^J .
\end{equation}

考虑 $M$ 是带圆度规的球面 $S^N$ 的情形。

仿照我们在课上对同一靶空间的玻色模型所作的步骤，确定本模型的大 $N$ 行为。

% [Tutor Note: 记号上有两处须留意。其一，原文的 $\mathbb{R}^{2,2}$ 应是 $\mathbb{R}^{2|2}$（两个玻色坐标加两个费米坐标的超空间）之排印，全书他处均按 $2|2$ 理解；若读成号差为 $(2,2)$ 的四维时空，$d^2x\,d^2\theta$ 的测度就讲不通。其二，$I,J$ 是靶空间 $M$ 的指标，$D_\pm$ 是超空间的协变导数（不是时空导数），$d^2\theta$ 积分给出 $D$ 项；$g_{IJ}$ 必须对称，且只有当 $X$ 是超场、度量为 $X$ 的函数时才能把 $g_{IJ}$ 提到导数外面。把 $I,J$ 误当时空指标，会写出 $D_+X_\mu$ 这类不存在的对象。]

% ===== 原文 p.287 =====

提示：以明显超对称的方式把 Lagrange 乘子纳入超空间。但到做路径积分时，先做 theta 积分，化为普通的拉格朗日量与路径积分，然后按至今应当已经熟悉的方式推进。（我们还没有发展出在超空间中计算 Feynman 图的超图技巧；而且对像本题这样简单的大多数问题，这类技巧反正也不会带来多少简化。）

这一显式程序应当让人看清：我们研究过的两个看似相似的机制（在玻色 $\sigma$ 模型与 Gross–Neveu 模型中）实际上是由超对称相互联系着的。

\subsection*{Problem 4}

证明前一题所考虑的模型尽管有质量间隙，却具有一个守恒荷，其类型是 Coleman–Mandula 定理所「不允许」的，即它不能写成 Poincaré 李代数的一个元素与一个与 Poincaré 对易的算符的线性组合。

\subsection*{Problem 5}

在靶空间为球面 $S^n$ 的二维 $\sigma$ 模型中，对大的 $n$，为了量子化该理论，我们引入了一个 Lagrange 乘子场 $\sigma$。

(a) 利用运动方程，用原始变量把 $\sigma$ 确定为一个算符。

(b) $\sigma$ 场的量纲是什么？因此你预期 $\langle\sigma(q)\sigma(-q)\rangle$ 在大 $q$（$q$ 为欧氏动量）下的行为如何？与实际行为比较。

% [Tutor Note: (a) 的做法值得写清：把约束写成 $L=\frac{1}{2g^2}\int\big[(\partial_\mu X)^2+\sigma(X^2-1)\big]$，对 $X^I$ 变分得 $-\partial^2X^I+\sigma X^I=0$，与 $X^2=1$ 联立即得 $\sigma=X\cdot\partial^2X=-\partial_\mu X\cdot\partial_\mu X$（最后一步用了 $X\cdot\partial_\mu X=0$）。所以 $\sigma$ 是复合算符 $-(\partial X)^2$，其期望值正比于能量动量张量的迹；这也解释了 (b) 中「按经典量纲 $\sigma$ 为 $2$」的来历。至于大 $q$ 行为，本题要你对比的正是：朴素量纲分析给出的幂次与真实幂次不符，差额即 $\sigma$ 的反常维数——不要满足于只写出 $q^{-4}$ 的量纲估计。]

\subsection*{Problem 6}

\footnote{本题大部分内容（最后一段除外）的参考文献是 S. Coleman, \textit{More about the massive Schwinger model}, Annals of Physics \textbf{101} (1976), 239--267。}%
在二维中考虑一个 $U(1)$ 规范理论，其中有\emph{两}种费米子 $\psi_i$，$i=1,2$，处于 $U(1)$ 的同一表示，裸质量为 $m_i$。于是拉格朗日量为

\begin{equation}
  \tag{5}
  \label{eq:ex-5}
  L_\psi(\psi_1,m_1)+L_\psi(\psi_2,m_2)+\int\left(\frac{F^2}{4e^2}+i\frac{F}{2\pi}\right)
\end{equation}

其中 $L_\psi$ 就是我们在课上考虑过的拉格朗日量：

\begin{equation}
  \tag{6}
  \label{eq:ex-6}
  L_\psi(\psi,m)=\int\left(\overline{\psi}_+D_-\psi_++\overline{\psi}_-D_+\psi_-+m\overline{\psi}_-\psi_++\overline{m}\,\overline{\psi}_+\psi_-\right)d^2x .
\end{equation}

% [Tutor Note: (5) 的拓扑项印作 $i\frac{F}{2\pi}$，其中并无显式的 $\theta$；但 (a)(b) 通篇讨论「$\theta$ 依赖」与「$\theta$ 的有效值移动」，可见这里应读作 $i\theta\frac{F}{2\pi}$（二维中也可把 $\theta$ 通过手征转动吸收进复质量的相位 $\arg m_i$，此时 (6) 里 $m$ 与 $\overline m$ 分开写就并非多余）。存疑处仅此一处，公式本体仍照原样录出。另注意 (5) 只在第二项写了积分号而省略了 $d^2x$，(6) 则把 $d^2x$ 写在末尾，二者是同一约定，不要以为是笔误。]

(a) 证明：只要有一个 $m_i$ 为零，$\theta$ 依赖就消失。

(b) 证明：在 $m_1$ 变得很大而 $m_2$ 固定的极限下（或者让 $m_2$ 趋于零而 $m_1$ 变得很大），理论约化为我们在课上研究过的那个理论，但 $\theta$ 的有效值可能发生移动。

(c) 定性地描述两个 $m_i$ 都很大时的行为。

(d) 描述两个 $m_i$ 都很小时的行为。

(e) 证明：若 $m_i$ 相等，理论具有 $SU(2)$ 对称性。现在，通过考虑 $m_i$ 相等且很小的情况，推断：二维 sine–Gordon 模型，即具有势
\[
  V(\phi)=\mu^2\cdot\cos(\beta\phi)
\]
的标量场模型，在某个 $\beta$ 值上具有 $SU(2)$ 对称性；你应当把这个 $\beta$ 值找出来。

% ===== 原文 p.288 =====

(f) 取 $\mu$ 趋于零的极限。此时 (e) 中提到的 sine–Gordon 理论变成无质量、共形不变的自由场论；由此推断：正确半径（等价地：固定半径 $2\pi$，并取动能项的正确系数 $k$，即 $L=k\int d^2x\,|d\phi|^2$）的二维自由无质量标量理论具有 $SU(2)$ 对称性。由于共形不变性，这个理论实际上表示的是 $SU(2)$ 的仿射化（affinization）（更确切地说，是它的两套拷贝，一套给左行波、一套给右行波），而不只是有限维群 $SU(2)$。这被称为 $SU(2)$ 仿射李群的顶点构造（vertex construction）。

要得到它，我们其实不必绕道规范理论。更一般地，只需考虑若干无质量自由费米子，对它们应用 Bose–Fermi 对应（Bose–Fermi correspondence），就能得到 $SO(2n)$（以及稍加修改后 $SU(n)$）的顶点构造。

\subsection*{Problem 7}

在二维中考虑一个 $SO(n)$ 规范理论，其标量场 $\phi$ 由 $n-1$ 份 $SO(n)$ 的 $n$ 维表示构成，势能为

\begin{equation}
  \tag{7}
  \label{eq:ex-7}
  V=\sum_{i=1}^{n-1}\lambda_i\left[(\phi_i,\phi_i)-a_i^2\right]^2+\sum_{i<j}c_{ij}(\phi_i,\phi_j)^2 .
\end{equation}

$\lambda_i$ 与 $c_{ij}$ 全都是正的。这里 $(\ ,\ )$ 是 $SO(n)$ 的 $n$ 维表示上的不变二次型。

由于 $\pi_1(SO(n))=\mathbb{Z}/2\mathbb{Z}$，该理论有一个双值的「$\theta$ 角」。请解释：在弱耦合下，你将如何计算 $\theta$ 取两个可能值时的两个真空之间的能量差。\footnote{如果你不确定当 $\pi_1(G)$ 是有限群时该如何理解 theta，可以查阅 E. Witten, \textit{Theta vacua in two-dimensional quantum chromodynamics}, Nuovo Cimento A \textbf{51} (1979), 325--338。}

% [Tutor Note: 两个易错点。第一，$\pi_1(SO(n))=\mathbb{Z}/2\mathbb{Z}$ 只对 $n\ge3$ 成立；$n=2$ 时 $SO(2)\cong U(1)$，$\pi_1=\mathbb{Z}$，$\theta$ 不是双值的——原文默认了 $n\ge3$。第二，(7) 中 $\lambda_i$ 与 $a_i$ 用同一个指标 $i$ 配对（每个 $SO(n)$ 矢量各有自己的尺度 $a_i$，故 $n-1$ 个矢量并不等价），而交叉项只对 $i<j$ 求和；若顺手写成 $i\ne j$，每个交叉项会被算两次，整个势就多出一个因子 $2$，真空流形（应是若干 $S^{n-1}$ 的乘积）随之变形。]

\subsection*{Problem 8}

\footnote{本题的参考文献是 E. Witten, \textit{Some properties of the $(\overline{\psi}\psi)^2$ model in two dimensions}, Nucl. Phys. B \textbf{142} (1978), 285。}%
在二维中考虑 $(\psi_+\psi_-)^2$ 模型（在 [I-Renorm Groups] 的第 4 讲中引入），其中的场是二维费米子 $\psi_i$，拉格朗日量为

\begin{equation}
  \tag{8}
  \label{eq:ex-8}
  L=\int d^2x\left(\sum_i\left(\psi_{+i}\partial_-\psi_{+i}+\psi_{-i}\partial_+\psi_{-i}\right)+g\sum_i\left(\psi_{+i}\psi_{-i}\right)^2\right).
\end{equation}

% [Tutor Note: 这里的 $\psi_{\pm i}$ 是不带 bar 的 Weyl 分量（本模型是手征的，$\overline\psi$ 不出现），$\partial_\pm=\partial_0\pm\partial_1$；若按四维习惯给 $\psi$ 加 bar，或把 $\partial_+$ 与 $\partial_-$ 对调，$S=\psi_{+1}\cdots\psi_{+n}$ 的守恒律立刻变成平凡或错误。另注意 $S$ 是 $n$ 个 Grassmann 数之积，$n$ 为奇数时它的整体符号依赖于因子的排序约定，取定一种排序后全篇不得再换——这一点在做 (b)（$n=4$）与 (d)（$n=3$）时会直接影响 $[K,R]$ 的符号。]

(a) 考虑经典守恒律 $\partial_-S=0$，其中 $S=\psi_{+1}\psi_{+2}\cdots\psi_{+n}$。通过考察它在量子力学中的变形，证明该模型是可积的（integrable）。

这样一来，你就得到一个守恒荷 $R$，满足 $[K,R]=(n/2-1)R$，其中 $K$ 是 Lorentz 群的生成元（二维 Poincaré 群的正规子群 $SO(1,1)$）。把 $+$ 与 $-$ 互换，你又得到另一个守恒荷 $R'$，满足 $[K,R']=-(n/2-1)R'$。

(b) 现在取 $n=4$，于是 $(n/2)-1=1$。因此 $R$ 与 $R'$ 具有与动量分量 $P_+$ 和 $P_-$ 相同的 Lorentz 变换律。

% ===== 原文 p.289 =====

你应当能把上面 (a) 的论证加强，以证明当 $n=4$ 时理论有两个能量动量张量 $T_{mn}$ 与 $T'_{mn}$：其中 $T$ 是标准的能量动量张量，其适当的积分给出 Poincaré 生成元；而 $T'_{mn}$ 的性质是，它的适当矩（moment）给出 $R$ 与 $R'$。但现在，把 $T'$ 与 $SO(1,1)$ 的生成元「缩并」，你就得到一个新的「Lorentz 生成元」$K'$；证明通常的 Poincaré 算符连同 $R$、$R'$ 与 $K'$ 生成两套相互对易的 Poincaré 群拷贝。（要得到对易的拷贝，你必须取「旧」算符与「新」算符的和与差。）

证明 $P_+$ 与 $R$（或 $P_-$ 与 $R'$）的任何线性组合都不为零，从而每一套 Poincaré 群拷贝都非平凡地作用。事实上，你应当能看出：$SO(4)$ 的外自同构（outer automorphism）（在本理论中它作用为 $\psi_1\to-\psi_1$，而对 $i>1$ 有 $\psi_i\to\psi_i$）交换这两套 Poincaré 拷贝。

由此推断：当 $n=4$ 时，本理论作为两个退耦理论的乘积是「可约的」——在其中一个理论里一套 Poincaré 起作用，在另一个里则是第二套。

% [Tutor Note: $[K,R]=(n/2-1)R$ 中的 $-1$ 常被漏掉。$S$ 本身的标度维数是 $n/2$（$n$ 个 $\psi_+$ 各贡献 $1/2$），故 $[K,S]=(n/2)S$；但 $R=\int dx^1S$ 还要对空间坐标积分，而 $dx^1$ 在 $SO(1,1)$ 下带权 $-1$，于是 $R$ 的维数是 $n/2-1$。若只算 $S$ 的维数，会得到 $[K,R]=(n/2)R$，那么 (b) 中 $n=2$ 时 $R$ 就该像 $P_+$ 一样是动量——与已知的 $(\psi_+\psi_-)^2$ 模型不符。顺带一提，$K$ 是反厄米的（boost 生成元），故 $R$ 与 $R'$ 的「同律」只到实系数为止。]

\subsection*{Problem 9}

课上我们在二维中考虑过这样一个理论：含 $n$ 个手征超场（chiral superfield）$A_i$，$i=1,\dots,n$，超势（superpotential）$F(A_1,\dots,A_n)$ 是 $n$ 次齐次函数。然后我们构造了该理论的一个 $\mathbb{Z}_n$ orbifold，其中对作用为 $A_i\to e^{2\pi i/n}A_i$ 的 $\mathbb{Z}_n$ 群作商。该理论具有左行与右行的 $R$ 对称群 $U(1)\times U(1)$，你对此应当熟悉。于是 Hilbert 空间是双分次的（bigraded）。

把该理论表述在圆 $S^1$ 上（对所有场取超对称的、周期的边界条件）。理论有一个有限维的超对称基态空间；事实上，出于课上应当已经清楚的理由，这个空间与方程 $F=0$ 在 $\mathbb{P}^{n-1}$ 中所定义射影超曲面的上同调（cohomology）密切相关。两者都是双分次的：一侧的双分次来自 $R$ 对称性，另一侧的双分次来自超曲面上同调的 Hodge 分解；两个双分次将（在平移 $n/2$ 之后）吻合，即：通常所谓的 $H^{p,q}$ 映到 $R$ 荷为 $p',q'=p-n/2,\;q-n/2$ 的态。

计算 orbifold 理论的超对称基态空间。这有不同的部分：

(1) 有 $n-1$ 个 twisted sector（扭化扇区）。每个 twisted sector 贡献多少个零能态？为求出 twisted sector 中基态的 $p',q'$ 值，你可以注意以下几点：(a) 由于没有零模，微扰论是直截了当的；(b) 计算 $p'$ 与 $q'$ 时可以取弱耦合极限；(c) 因此可以用自由场论；(d) 于是该计算与超弦理论中比较 Ramond 扇区与 Neveu–Schwarz 扇区时出现的计算极为相似；(e) 你本质上只需计算一个相当平凡的一维算符的 eta 不变量。

% ===== 原文 p.290 =====

(2) 在 untwisted sector（非扭化扇区）中，微扰论并不直截了当。不过在这种情况下，非零模不作任何重要贡献（不像前一种情形有 eta 不变量），因此，把 $S^1$ 想得很小、丢掉非零模、约化到 $0+1$ 维，并不会丢失任何本质的东西。于是我们处在经典数学的领域。用你所拥有的量子力学、$L^2$ 上同调、奇点理论或别的什么知识，确定这个 $0+1$ 维问题的零能态谱。

现在把 (1) 与 (2) 的答案合并起来，并与射影超曲面的上同调比较。

下面两道习题涉及二维 $\mathcal{N}=2$ 超对称理论的扭曲（twisted）版本。回忆：在这样的扭曲理论中，人们只考虑与某个 $Q$ 对易的算符，而 $Q=Q_++Q_-$ 或 $Q=\overline{Q}_++Q_-$（或其中之一的反厄米共轭）。若在平坦的 $\mathbb{R}^2$ 上工作，则 $Q$ 与 $S+S'$ 对易，其中 $S$ 是 $\mathbb{R}^2$ 的转动生成元，$S'$ 是左右两套 $R$ 对称性之和或差。

\subsection*{Problem 10}

考虑含手征超场 $\Phi_i$、具有某个多项式超势 $W(\Phi_i)$ 的理论。我们假定 $W$ 只有有限多个临界点（可以是退化的）。拉格朗日量为

\begin{equation}
  \tag{9}
  \label{eq:ex-9}
  L=\int d^2x\,d^4\theta\;\overline{\Phi}_i\Phi_i+\left(\int d^2x\,d^2\theta\,W(\Phi)+\mathrm{c.c.}\right).
\end{equation}

% [Tutor Note: (9) 的第一项是 D 项（$d^4\theta$ 取出 $\theta^4$ 分量），第二项是 F 项（$d^2\theta$ 取出 $\theta^2$ 分量），二者量纲之和才使 $L$ 无量纲——写漏 $\overline{\Phi}_i\Phi_i$ 上的指标缩并（此处取平坦 Kähler 势 $\sum_i|\Phi_i|^2$，一般 Kähler 流形应换成 $K(\Phi,\overline\Phi)$）或把 $d^4\theta$ 误写成 $d^2\theta$，都会使 $W$ 的临界点数目与算符维数全部错位。「$W$ 只有有限多个临界点」是后面把 (10) 约化为有限维积分的前提：它保证手征环 $R=\mathbb{C}[\Phi]/\partial W$ 是有限维的。]

我们回顾：该理论可由四维理论作量纲约化（dimensional reduction）得到，四维坐标为 $x^1,x^2,\dots,x^4$；约化到二维的办法是让场与 $x^3,x^4$ 无关。设 $S$ 是 $1-2$ 平面的转动生成元，$S'$ 是 $3-4$ 平面的转动生成元。对任何 $W$，$S'$ 都是对称性。用 $S'$ 作扭曲，得到拓扑场论（topological field theory），其中 $Q$ 不变的算符是手征超场（加上它们在 $q=1,2$ 时的 $q$-形式后代（descendant））；由于 $S'$ 对任何 $W$ 都是对称性，这个理论对任何 $W$ 都有定义。

令 $\mathcal{O}_i$ 为手征超场（注意可以只取 $\Phi_i$ 的多项式函数，不含导数），在亏格零时计算路径积分

\begin{equation}
  \tag{10}
  \label{eq:ex-10}
  \left\langle\prod_{i=1}^{s}\mathcal{O}_i(x_i)\right\rangle .
\end{equation}

作为预备性的约化，注意利用手征环（chiral ring）可以约化到 $s=1$ 的情形。

要计算亏格零的路径积分，主要步骤是丢掉所有场的非零模，从而约化为一个有限维积分。如果你不知道如何做、或如何解释所得的有限维积分，请先用只有一个手征超场的情形练手。

\subsection*{Problem 11}

首先考虑二维中规范群为 $G$ 的 $\mathcal{N}=2$ 规范理论。其场强是 twisted 手征超场 $\Sigma=\sigma+\dots$。

% ===== 原文 p.291 =====

这个理论也可以由四维作量纲约化而来。四维中存在一个 $R$ 对称性 $S''$。约化到二维后，我们得到一个 twisted 拓扑理论，其中转动由 $S+S''$ 生成，而拓扑对称性是 $\overline{Q}_++Q_-$；该理论中的可观测量是 twisted 手征超场（而不是像前一题那样是手征超场）。

我们证明了以下几点：

(a) $\sigma$ 的不变多项式（以及它们的 $q$-形式后代）给出拓扑可观测量。

(b) 这样的多项式 $P$ 对应于平坦联络（flat connection）模空间（moduli space）上的一个上同调类 $w_P$。

(c) 关联函数（correlation function）由这类上同调类之积在模空间上的积分给出。（考虑 $q$-形式后代很重要，否则一切都为零；这一论断最初由 Newstead 对 $G=SU(2)$ 猜测。）

现在加入某些 $G$ 表示中的手征超场 $\Phi_i$ 来修改它。注意（与问题 (1) 中我们所作的、允许任意超势的扭曲相反）这里我们作的是 $S+S''$ 扭曲，它禁止一般的超势。\footnote{若想包含超势，则 (i) 它在 $\Phi_i$ 的某个 $\mathbb{C}^*$ 作用下必须具有权 $2$；(ii) 必须把该 $\mathbb{C}^*$ 的生成元加进 $S''$，从而改变几何 $R$ 对称性在手征超场上的作用方式。}因此，取超势为零。

% [Tutor Note: 原文此处写「与问题 (1) 中……相反」，但习题编号是从 1 到 11 的，此处所指显然是前一道题 Problem 10（它的扭变由 $S'$ 给出、允许任意超势）。这应是排印笔误（把「10」误作「(1)」），译文照录并在此说明。另注意本题与前一题的扭变生成元不同：Problem 10 用 $S'$（3-4 平面转动），故 $Q=Q_++Q_-$、可观测量是手征超场；本题用 $S''$（四维 $R$ 对称性），故拓扑对称性是 $\overline{Q}_++Q_-$、可观测量是 twisted 手征超场，两者不可混用。]

(a) 与 (b) 没有变化。然而 (c) 会被修改。在含手征超场的扭曲理论中，算符 $P^{(q)}$（在 $\Sigma$ 中的 $q$-闭链上积分）的关联函数，将是同一个上同调类之积在平坦联络模空间上的积分，再乘上一个依赖于手征超场的新类。你能认出这个类吗？
